\chapter{Methodology}
\label{c:methodology}

\noindent
The research questions and hypotheses are introduced first, followed by the corpus, the typology
of ambiguity types in texts, and the author-framing manipulation. The chapter ends with the
experimental design as it was actually carried out and an account of how that design differs
from what was originally proposed.

\section{Research Questions}
\label{meth:rqs}

\noindent
The thesis addresses three research questions.

\begin{itemize}
    \item \textbf{RQ1 (Framing Effect in Humans).} Does perceived author identity significantly
    influence human emotion judgments of ambiguous texts?

    \item \textbf{RQ2 (Framing Effect in LLMs).} Does perceived author identity significantly
    influence emotion classifications produced by large language models?

    \item \textbf{RQ3 (Human--Model Alignment).} Are the framing effects of LLMs similar in
    direction and magnitude to those of human annotators?
\end{itemize}

\section{Research Hypotheses}
\label{meth:hypotheses}

\noindent
These questions are operationalised in four hypotheses.

\begin{itemize}
    \item \textbf{H1 (Human Framing Hypothesis).} Human emotion judgments will differ significantly
    across author-framing conditions for ambiguous texts, but not (or substantially less) for
    unambiguous control texts.

    \item \textbf{H2 (Model Framing Hypothesis).} LLM emotion classifications will shift across
    author-framing conditions, reflecting sensitivity to minimal author cues.

    \item \textbf{H3 (Ambiguity Sensitivity Hypothesis).} Framing effects will be stronger for
    author-relevant ambiguous texts than for author-independent ambiguous texts.

    \item \textbf{H4 (Alignment Hypothesis).} The direction of framing-induced shifts in LLM
    classifications will partially align with human judgment shifts, but the magnitude of these
    shifts will differ.
\end{itemize}

\noindent
The scoring of the evidence for H4 is affected by two aspects of its wording. First, H4
predicts \emph{partial} directional alignment, so a positive but weak directional correlation is
counted \emph{for} the hypothesis, rather than against it. Second, H4 predicts that magnitudes
will \emph{differ}. Human and model shift sizes that turned out similar would thus be a surprise
with respect to the hypothesis as stated, not an affirmation of it. Chapter~\ref{c:result} scores
the two components individually.

The approved proposal contains no numbered hypothesis about prompt structure. Prompt-configuration
effects are thus treated as an exploratory question, not as a test of a formal hypothesis.

\section{Corpus Selection and Dataset Construction}
\label{meth:corpus}

\noindent
The source of the data is the crowd-enVent corpus \cite{Troiano2023CrowdEnVent}, which comprises
short self-authored descriptions of emotional events. Every text has two layers of annotation:
one label provided by the author of the text, indicating the emotion they were trying to express,
and labels provided by five independent readers who did not see the author's label. These two
layers make the corpus suitable for this study: disagreement can be assessed both between author
and readers and among the readers themselves.

The corpus uses eleven categories of emotion, consisting of \emph{anger, disgust, fear, guilt,
joy, pride, relief, sadness, shame, surprise} and \emph{trust}. All eleven are kept. The set is
small enough for human annotators and models to work with, and deleting any category could
suppress exactly the disagreement patterns the study aims to observe. This fixed, alphabetically
ordered set of eleven is the dimension set for every probability vector built in
Chapter~\ref{c:implementation}.

To operationalise ambiguity, two quantities were calculated per text. \textbf{Author--reader
agreement} asks whether the reader majority recovers the emotion the author claims to have
intended; a discrepancy means the text failed to convey what its author intended. \textbf{Reader
consensus} (\emph{majority share}) is the fraction of the five readers who selected the majority
emotion. Since there are exactly five readers, this measure only takes the values $0.2$, $0.4$,
$0.6$, $0.8$ and $1.0$, a coarseness that directly constrains where a threshold can sit.

Reader consensus does not cluster at the top of that scale. In the candidate corpus, a majority
share of $0.6$ (three readers of five) occurs most frequently. Disagreement is the typical case
and not an edge case, and this is the empirical basis for the distributional representation
adopted in Section~\ref{impl:vectors}.

\subsection{The classification rule}
\label{meth:rule}

\noindent
Three decisions are applied in a fixed order. Their thresholds and ordering are encoded in
\texttt{scripts/classify\_ambiguity\_pools.py}. \textbf{The persona-flip input for each text (Step
2 below) was generated by an LLM judge, not by human pilot annotators}: the judge was shown the
same text under two candidate personas and asked for the emotion it read into each, and the two
judgments were compared directly. This LLM-judging step is described in the dataset-construction
notebook, \texttt{complete\_analysis.ipynb}; the specific model used as the judge is not recorded
in the materials retained for this thesis and could not be reconstructed at the time of writing.
The full derivation, including a threshold sensitivity analysis, is documented in
\texttt{00\_corpus\_and\_sampling.ipynb}.

\textbf{Step 0: exclude off-label texts.} Texts whose author emotion or reader majority is
\emph{boredom} or \emph{no-emotion} are dropped, since these fall outside the eleven-emotion set.
Forced into the taxonomy, they produce persona-to-persona differences that are a labelling
artefact, not a real interpretive shift. The step removes 215 of the 1{,}200 candidates and
leaves 985 eligible.

\textbf{Step 1: threshold first.} A text is \textbf{unambiguous} if
$\text{majority share} \geq 0.8$ \emph{and} the readers recover the author's own label. This is
decided purely from the original corpus's reader data.

\textbf{Step 2: flip test, among the remainder only.} Of the texts not classified as
unambiguous, those whose two candidate personas yield different judged emotions from the LLM judge
are \textbf{author-relevant ambiguous}; the rest are \textbf{author-independent ambiguous}.

\noindent
\textbf{The ordering is deliberate.} A persona flip does \emph{not} override a strong,
author-matching reader consensus: a text the readers already agree on remains unambiguous even
when the LLM judge is moved by a persona flip. With the order reversed, the flip test could shift
high-consensus texts into the author-relevant category, and that category's definition would then
become circular with respect to H3.

\textbf{A filtering inconsistency at this step.} Step 0 filters out texts whose author emotion or
reader majority does not lie within the eleven-emotion set, on the grounds that forcing an
off-label judgment into the taxonomy is an artefact of labelling and not a true interpretive
shift. The same criterion was not applied to the LLM judge's own answers at Step 2: of the 100
texts classified as author-relevant, 24 have the judge answering \emph{no-emotion} for at least
one persona, precisely the type of off-label judgment Step 0 treats as unreliable elsewhere. These
texts were not re-screened against that criterion, so the author-relevant category is filtered
less consistently than the unambiguous category.

With a five-valued measure, only three threshold settings are defensible, and $0.8$ is one of
them. At $0.6$ the unambiguous pool would take in texts on which only three of five readers agreed,
weakening the control category the design relies on; at $1.0$ it would admit only unanimous
texts. Every one of the three leaves more than 100 candidates per pool, so availability does not
force $0.8$; it was chosen as the stricter of the two middle options.

\subsection{Drawing the sample}
\label{meth:sampling}

\noindent
The sample was not a single fresh draw and should not be described as one. It was
produced by \textbf{reconciliation} against an existing sample: the current sample is kept, every
text in it is reclassified under the rule above, any pool of more than 100 texts is pruned back to
100 using the fixed seed ($42$), and shortfalls are filled by drawing from eligible candidates not
already used. Three texts set aside as one-shot and few-shot prompt examples are excluded
throughout, so that no sampled text can appear inside its own prompt. This is asserted to result
in 300 texts, 100 per pool, with no duplicates.

Because the pools differ in size, so do the sampling fractions. The design puts equal $n$ per
category ahead of proportional representation, so that category comparisons are not confounded
with unequal sample size.

\textbf{What is and is not reproducible.} The classification rule is fully reproducible and is
verified against all 300 sampled texts. The \emph{specific membership} of the sample is not
recoverable from the seed alone, because it depends on the prior sample's contents and on the
candidate ordering at the time each replacement was drawn. \texttt{sampled\_texts.csv} is
therefore the authoritative record of which texts were used.

Seven verification checks on the sample (size, balance, absence of reserved examples, absence
of duplicates, pool eligibility, agreement with the classification rule, and absence of off-label
texts) all pass and are enforced as assertions rather than advisory output in
\texttt{00\_corpus\_and\_sampling.ipynb}.

\textbf{One limit of the off-label filter.} Step 0 screens two labels per text, the author's
own and the selected reader-majority label, not the individual reader votes. A text can
therefore enter the sample while still containing reader votes outside the eleven-emotion set, and
27 such votes survive across 25 texts, so \textbf{275 of the 300 sampled texts retain all five
in-inventory reader votes}. Those votes are excluded when the original-corpus vectors are built,
which is why vector construction tracks $n$ per vector rather than assuming five throughout
(Section~\ref{impl:vecdef}).

\section{Ambiguity Types and Author Relevance}
\label{meth:ambiguity}

\noindent
Not all ambiguous texts are ambiguous for the same reason. Some are uncertain because
author-related context is missing; others describe an underspecified or emotionally complicated
situation in which the author's identity is not material. H3 relies on this distinction, since it
predicts that framing effects concentrate in the former.

Texts were accordingly sorted into three groups, used consistently throughout the thesis:

\begin{itemize}
    \item \textbf{Author-relevant ambiguous texts}, where information about the author's role,
    perspective, or position within society could plausibly shift the perceived emotion. A
    description of a risky action may read as pride from the agent who took it and as fear or
    worry from a spectator.

    \item \textbf{Author-independent ambiguous texts}, where several emotions remain plausible
    regardless of who wrote the text, typically attributable to a lack of situational detail
    rather than to an authorship issue.

    \item \textbf{Unambiguous control texts}, drawn from the high-agreement pool, where the writer
    and the readers agree strongly on one emotion.
\end{itemize}

\noindent
The final sample consists of \textbf{300 texts, balanced at exactly 100 per category}. This
balance is a property of the sampling design and is verified programmatically as part of the
quality-control gate described in Section~\ref{impl:qc}.

One property of this categorisation constrains what H3 can claim, and it must be stated
plainly. Texts were assigned to the author-relevant group partly on the basis of
an LLM judge's own sensitivity to whether author framing would invert the reading: the corpus
preparation records, for each text, two candidate persona readings from that judge and whether
they differ. Texts were therefore \emph{selected} into that group partly on the property H3
subsequently measures. H3 is consequently a \textbf{transfer comparison}: it asks whether a
grouping based on the judgments of one LLM predicts framing effects in newly collected human
data. That author-relevance is the operative principle is therefore not an independent
discovery, and Chapter~\ref{c:result} does not present it as one.

The impact is sharper on the model side than on the human side. Because the category itself was
defined by whether an LLM's judgment flips between personas, the finding that the \emph{model
ensemble's} framing-effect magnitude is ordered author-relevant $>$ author-independent $>$
unambiguous (Section~\ref{res:h2magnitude}) is partly true by construction: texts were selected into the
author-relevant pool precisely because at least one LLM's reading was sensitive to the persona
pair. The ordering holds across all six prompt configurations, and for an ensemble of seven
distinct models rather than just the one used as classification judge, which suggests that the
pattern is more than an artefact of the selection. Nonetheless, the model-side ordering claim must
be read as partially circular rather than as an independent confirmation, just as the
human-side H3 result is treated as a transfer comparison rather than a fresh discovery.

\section{Author-Framing Manipulation}
\label{meth:framing}

\noindent
Perceived authorship is the independent variable. The content of every text is held
\emph{exactly} constant; only a short author cue attached to it varies, so no difference in
emotion judgment can come from a change in the text itself. The cue is appended after the
text rather than placed before it (Section~\ref{impl:stimulus} gives the exact format).
Which side of the text the cue sits on is not expected to matter for the manipulation's logic,
but the earlier description of it as \emph{preceding} the text was not consistent with how it was
actually presented. Nor does the manipulation by itself isolate identity from every other kind of information
the cue carries: a persona description can also supply substantive situational detail beyond who
the author is (for instance, a persona contrasting someone who has been struggling all year with
someone expecting to pass easily differs in circumstance as well as in identity), so the design
measures sensitivity to the provided author/persona framing as a whole, not a pure identity
mechanism in isolation.

Four framing conditions were used:

\begin{itemize}
    \item \textbf{Neutral}: no author information is provided.
    \item \textbf{Generic human}: the text is attributed to an unspecified human author.
    \item \textbf{Persona A} and \textbf{Persona B}: the text is attributed to a person with a
    defined social role.
\end{itemize}

\noindent
The generic human condition is more than an extra level; the analysis depends on it. Neutral and
persona prompts differ in more than persona content, since persona prompts include an author
line that is absent from neutral prompts. A neutral-to-persona comparison therefore confounds
\emph{who the author is} with \emph{whether an author is named at all}. The generic human
condition separates the two, and Section~\ref{res:generic} uses it for exactly that purpose.

Persona selection was not straightforward. There is no single fixed pair, such as \emph{politician} and
\emph{comedian}, that applies equally to every text; a domestic anecdote has no natural connection
to either role. Personas were therefore assigned per text to suit its subject matter, with the two
personas of a pair chosen to have contrasting interpretive expectations (one related to
accountability, one to latitude, for instance), so that a framing effect, if present, has a
reasonable chance of emerging. Both personas for every text, along with the reasoning behind the
pairing, are recorded explicitly in the dataset.

\section{Realised Experimental Design}
\label{meth:design}

\noindent
Table~\ref{tab:design} summarises the design as executed.

\begin{table}[h]
\centering
\footnotesize
\caption{Realised experimental design.}
\label{tab:design}
\begin{tabular}{lp{7.5cm}l}
\hline
\textbf{Factor} & \textbf{Levels} & \textbf{Count} \\
\hline
Texts & balanced across 3 ambiguity categories & 300 (100 per category) \\
Emotion labels & fixed, alphabetically ordered & 11 \\
Framing (models) & neutral, generic human, persona A, persona B & 4 \\
Framing (humans) & neutral, persona A, persona B & 3 \\
Models & 7 (proprietary, commercial and open-weight) & 7 \\
Prompt variants & zero-shot, one-shot, few-shot & 3 \\
Label formats & numbered list, inline & 2 \\
\hline
Model judgments & $300 \times 4 \times 3 \times 2 = 7{,}200$ cells, each an ensemble vector aggregating up to 7 models' raw answers (50{,}400 raw answers total) & 7{,}200 ensemble vectors \\
Human judgments & $300 \times 3$ sides & 900 vectors \\
\hline
\end{tabular}
\end{table}

\noindent
\textbf{Human annotation.} Participants were recruited and presented with texts one at a time
together with the assigned author framing, selecting a label from the same fixed set of eleven
emotions used for the models. Texts and framings were distributed such that each text--framing
pair received several independent annotations from different people, while \textbf{no individual
annotator saw the same text more than once}. This property was specified in the approved
proposal, and the statistical analysis of the human data in Section~\ref{impl:stats} is based on
it: since each participant was shown only one framing of a given text, the conditions for a text
are rated by disjoint sets of individuals, and there is no within-text pairing across conditions
to preserve.

\textbf{Model evaluation.} Each model saw each text only once per framing condition and prompt
configuration, avoiding repeated exposure. The predicted label for each (text, framing, variant,
format, model) combination was recorded. Prompt configurations are always kept separate in any
analysis; they are only reported, or averaged, after being calculated independently.

\section{Deviations from the Approved Proposal}
\label{meth:deviations}

\noindent
The departures from the proposal are detailed here so that none is left implicit.

\textbf{Sample size.} The proposal expected 300--400 texts. The actual sample is 300, chosen to
allow an exact balance of 100 per category.

\textbf{Number of models.} The proposal outlined at least one proprietary, one commercial, and one
open-source model. Seven were used, which substantially supports RQ2: a framing effect that shows
up across seven independently developed systems is more difficult to attribute to a peculiarity
of a single model family.

\textbf{Number of prompt configurations.} The proposal specified at least two prompt variants.
Three variants were crossed with two label formats (six configurations), making it possible to
evaluate prompt sensitivity as an exploratory question in its own right.

\textbf{Ambiguity typing and persona mapping.} The proposal envisaged manual inspection of a
subset of texts, thematic clusters, and persona pairs assigned consistently within each cluster. The
realised classification instead employed a fixed consensus threshold and a documented pilot
persona-flip decision for the remaining texts; persona pairs were assigned text by text. The
category definition incorporates the pilot judgments, so the author-relevant category is partly
selected for the property tested by H3 (Section~\ref{meth:ambiguity}). Under this per-text
design, a cluster-to-persona mapping was not generated. The realised pairings and rationales
can be inspected only in the master sampling file; they are not catalogued in this manuscript.

\textbf{Asymmetry in framing conditions.} Models were evaluated under all four framings; humans
were annotated under three, omitting the generic human condition. The control analysis in
Section~\ref{res:generic} can therefore only be applied to models, and it is not possible to
provide a parallel human decomposition of author-presence versus persona-content.
Section~\ref{disc:limitations} returns to this limitation of the collected data.

\textbf{Inter-annotator agreement.} The proposal specified Fleiss' $\kappa$ \cite{Fleiss1971Kappa}.
This statistic does not require respondent identifiers: each item contributes its counts of
emotion votes. Since the realised number of votes per (text, framing) cell ranges from three to
seven, the conventional equal-rater calculation is restricted to the 754 cells with exactly three
votes. The other 146 cells are excluded from this descriptive calculation only, not from the main
distributional analysis. Kappa is calculated separately within each ambiguity pool and framing
condition across the eleven emotion labels (Section~\ref{res:agreement}). It corrects for the
group-specific marginal emotion frequencies, so comparing values between groups is a descriptive
rather than a causal claim about framing. The equal-three restriction can also select a
non-representative subset of persona cells.

\textbf{Model evaluation metric.} The proposal specified macro-F1 and percentage agreement
against human majority labels. Section~\ref{res:majoritybenchmark} provides both as a secondary
neutral-framing benchmark against the original corpus's reader majority, using only the 241
sampled texts with a strict majority among five readers. Each of the six prompt configurations is
scored individually and averaged by model; invalid labels are removed, and the count of retained
predictions is presented in the reproducible table. The benchmark answers a narrower question
than the main analysis: it uses a noisy reader majority as the gold label, filters out 59 texts
with no strict majority, and does not score persona responses against the separately
collected human votes. Model framing behaviour is instead characterised distributionally: by
agreement with human vectors at baseline (Section~\ref{res:vectorization}), and by the magnitude
and direction of the shift under framing (Sections~\ref{res:h2}, \ref{res:h4}).

\textbf{Qualitative case analysis.} The proposal outlined an analysis of individual strong-shift
and stable-text cases. One traceable example, selected for its observed human shifts and reported
in Section~\ref{res:example}, illustrates how the reading of the same text can change completely
under one persona and stay completely unchanged under another. This illustrative example is not a systematically
coded thematic analysis of cases and cannot be used to determine why the shift occurred.

\textbf{Explanations and corpus summaries.} All seven analysed models provided a nonempty
explanation for every one of their 7{,}200 responses (coverage is documented in
\texttt{results/\allowbreak model\_explanation\_coverage\_verified.csv}); one such response is
discussed in Section~\ref{res:example}. Neither an explanation-based thematic analysis nor a
model confidence analysis was conducted. Section~\ref{res:sampledescription} provides
author-label frequencies and text lengths for the final 300-text sample. These sample statistics
are not a proxy for a separately requested descriptive analysis of the full source corpus.

\textbf{Participant count.} The proposal projected 10--15 participants. The
realised study used a larger, less formally demarcated pool recruited across multiple collection
rounds (an early round and a later persona/neutral re-collection round, Section~\ref{impl:leak}),
with a target of three annotators per (text, framing) cell. Since respondent identifiers were not
stored (Section~\ref{impl:human}), the actual number of distinct participants cannot be
reconstructed after the fact, and no such claim is made. This is itself a limitation, revisited in
Section~\ref{disc:limitations}. The within-text permutation null of Section~\ref{impl:permutation}
does not require knowledge of the total number of people involved, but different people rating
the two conditions eliminates only the within-text pairing requirement; it is not, in itself,
evidence that the two conditions are exchangeable (Section~\ref{impl:permutation}).
